\matlabtitle{Исследование типовых динамических звеньев}

\section{Исследование типовых динамических звеньев}

В рамках данной лабораторной работы необходимо исследовать реальные объекты, сопоставим их с известными нам типовыми динамическими звеньями и получим их передаточные функции в стандартизированной форме (через коэффициенты усиления $K$ и постоянные времени $T$). Затем запишем аналитические выражения для временных (переходной и весовой) и частотных (АЧХ, ФЧХ, ЛАЧХ и ЛФЧХ) характеристик исследуемых звеньев. Проведем моделирование с целью получения графического представления всех перечисленных характеристик и сравним с теоретическими для данного типового звена.

\textbf{Вариант 27}

\subsection{ДПТ}

Даны уравнения двигателя постоянного тока независимого возбуждения:

$J\dot{\omega} =M$, $M=k_{\text{m}} I$, $I=\frac{U+\varepsilon_{\text{i}} }{R}$, $\varepsilon_i =-k_{\text{e}} \omega$.

Возьмем из \textbf{Таблицы 1} инструкции по выполнению данной л/р значения, которые соответствуют нашему варианту, для следующих величин:

• $k_{\text{m}} =0.3824$ Н·м/А — конструктивная постоянная по моменту;

• $k_{\text{e}} =0.3824$ В·с — конструктивная постоянная по ЭДС;

• $J=0.0019$ кг·м\textasciicircum{}2 — момент инерции ротора;

• $R=4.6296$ Ом — активное сопротивление обмоток ротора.

Входом объекта считаем $U(t)$, а выходом $\omega (t)$.

\textbf{Найдём передаточную функцию и определим тип звена:}

$$
J\dot{\omega} =k_m \cdot \frac{U-k_e \omega }{R}=\frac{k_m }{R}U-\frac{k_e k_m }{R}\omega ,~~J\dot{\omega} +\frac{k_e k_m }{R}\omega =\frac{k_m }{R}U\\
\frac{JR}{k_e k_m }\dot{\omega} +\omega =\frac{1}{k_e }U
$$

Передаточная функция: $W(s)=\frac{K}{Ts+1}$, где $K=\frac{1}{k_e }=\frac{1}{0.3824}=2.615,~~T=\frac{JR}{k_e k_m }=\frac{0.0019\cdot 4.6296}{0.3824\cdot 0.3824}=0.0602$.

Типовое звено - \textit{реальное усилительное}.

\textbf{Запишем аналитические выражения для временных и частотных характеристик:}

$$W(j\omega )=\frac{K}{T(j\omega )+1}=\frac{K-KT(j\omega )}{T^2 \omega^2 +1}=\frac{K}{T^2 \omega^2 +1}-\frac{KT\omega }{T^2 \omega^2 +1}j$$

Весовая: $w(t)={\mathcal{L}}^{-1} \lbrace W(s)\rbrace ={\mathcal{L}}^{-1} \left\lbrace \frac{K}{Ts+1}\right\rbrace ={\mathcal{L}}^{-1} \left\lbrace \frac{K}{T}\cdot \frac{1}{s+\frac{1}{T}}\right\rbrace =\frac{K}{T}e^{-\frac{t}{T}}$

Переходная: $h(t)={\mathcal{L}}^{-1} \left\lbrace \frac{W(s)}{s}\right\rbrace ={\mathcal{L}}^{-1} \left\lbrace \frac{K}{s(Ts+1)}\right\rbrace ={\mathcal{L}}^{-1} \left\lbrace \frac{K}{s}-\frac{K}{s+\frac{1}{T}}\right\rbrace =K\left(1-e^{-\frac{t}{T}} \right)$

АЧХ: 

$$
A(\omega )=|W(j\omega )|=\sqrt{{\text{Re}}^2 (W)+{\text{Im}}^2 (W)}=\sqrt{{\left(\frac{K}{T^2 \omega^2 +1}\right)}^2 +{\left(\frac{KT\omega }{T^2 \omega^2 +1}\right)}^2 }=
$$
$$
=\sqrt{\frac{K^2 +T^2 K^2 \omega^2 }{(T^2 \omega^2 +1)^2 }}=\frac{K}{\sqrt{T^2 \omega^2 +1}}
$$

ФЧХ:

$$\varphi (\omega )=\arg W(j\omega )=\text{atan2}(\text{Im}(W),\text{Re}(W))=\text{atan2}\left(-\frac{KT\omega }{1+T^2 \omega^2 },\frac{K}{1+T^2 \omega^2 }\right)$$


ЛАЧХ и ЛФЧХ находятся из формулы: $\lg (W(j\omega ))=\lg (A(\omega ))+j\lg (e)\phi (\omega )$.

ЛАЧХ: $L(\omega )=20\lg \left(A(\omega )\right)=20\lg K-20\lg \sqrt{T^2 \omega^2 +1}$

ЛФЧХ: $\varphi (\omega )=\text{atan}\left(\frac{\text{Im}(W)}{\text{Re}(W)}\right)=-\text{atan}(T\omega )$

Для моделирования реакции системы на step response будет использоваться схема на рис. \ref{stepr}, для impulse response - на рис. \ref{impr}.
(Для остальных объектов структурные схемы приводиться не будут, так как они полностью аналогичные, за исключением соответствующей ПФ).
\begin{figure}[H]
    \centering
    \includegraphics[width=1\linewidth]{/home/ns/Desktop/ITMO-courses/Linear_automatic_control_systems/lab_5/code/images/step.png}
    \caption{Стуктурная схема для step resp.}
    \label{stepr}
\end{figure}
\begin{figure}[H]
    \centering
    \includegraphics[width=1\linewidth]{/home/ns/Desktop/ITMO-courses/Linear_automatic_control_systems/lab_5/code/images/impulse.png}
    \caption{Стуктурная схема для impulse resp.}
    \label{impr}
\end{figure}
\textbf{Графики временных и частотных характеристик полученные в результате моделирования, их сравнение с построенными на основании аналитических выражений:}

Графики временных характеристик представлены на рис. \ref{time_dpt}.
\begin{figure}[H]
    \centering
    \includegraphics[width=1\linewidth]{/home/ns/Desktop/ITMO-courses/Linear_automatic_control_systems/lab_5/code/images/plot_dpt.png}
    \caption{Графики временных характеристик}
    \label{time_dpt}
\end{figure}

Графики частотных характеристик представлены на рис. \ref{freq_dpt}.
\begin{figure}[H]
    \centering
    \includegraphics[width=1\linewidth]{/home/ns/Desktop/ITMO-courses/Linear_automatic_control_systems/lab_5/code/images/freq_dpt.png}
    \caption{Графики частотных характеристик}
    \label{freq_dpt}
\end{figure}


\subsection{ДПТ 2.0}

Даны уравнения двигателя постоянного тока независимого возбуждения:

$J\dot{\omega} =M$, $M=k_{\text{m}} I$, $I=\frac{U+\varepsilon }{R}$, $\varepsilon =\varepsilon_{\text{i}} +\varepsilon_{\text{s}}$, $\varepsilon_{\text{i}} =-k_{\text{e}} \omega$, $\varepsilon_{\text{s}} =-L\dot{I}$.

Возьмем из \textbf{Таблицы 1} инструкции по выполнению данной л/р значения, которые соответствуют нашему варианту, для следующих величин:

• $k_{\text{m}} =0.3824$ Н·м/А — конструктивная постоянная по моменту;

• $k_{\text{e}} =0.3824$ В·с — конструктивная постоянная по ЭДС;

• $J=0.0019$ кг·м\textasciicircum{}2 — момент инерции ротора;

• $R=4.6296$ Ом — активное сопротивление обмоток ротора.

• $L=1.1194$ Гн — индуктивность обмоток ротора.

Входом объекта считаем $U(t)$, а выходом $\omega (t)$.

\textbf{Найдём передаточную функцию и определим тип звена:}

$$
J\dot{\omega} =k_m \cdot \frac{U+\varepsilon }{R}=k_m \cdot \frac{U-k_e \omega -L\dot{I} }{R},~~JR\dot{\omega} =k_m U-k_e k_m \omega -Lk_m \dot{I}
$$
$$
JL\ddot{\omega} +JR\dot{\omega} +k_e k_m \omega =k_m U,~~\frac{JL}{k_e k_m }\ddot{\omega} +\frac{JR}{k_e k_m }\dot{\omega} +\omega =\frac{1}{k_e }U
$$

Передаточная функция: $W(s)=\frac{K}{T^2 s^2 +2\xi Ts+1}$,\\
где $
K=\frac{1}{k_e }=2.615,~~T=\sqrt{\frac{JL}{k_e k_m }}=\sqrt{\frac{0.0019\cdot 1.1194}{0.3824\cdot 0.3824}}=\sqrt{0.01458}=0.1207,\\
\xi =\frac{R\sqrt{J}}{2\sqrt{Lk_e k_m }}=\frac{4.6296\cdot \sqrt{0.0019}}{2\cdot \sqrt{1.1194\cdot 0.3824\cdot 0.3824}}=\frac{0.2018}{2\cdot 0.3781}=0.2668.
$

Типовое звено - \textit{колебательное}.

\textbf{Запишем аналитические выражения для временных и частотных характеристик:}

$$W(j\omega )=\frac{K}{1-T^2 \omega^2 +j(2\xi T\omega )}=\frac{K(1-T^2 \omega^2 -j2\xi T\omega )}{(1-T^2 \omega^2 )^2 +(2\xi T\omega )^2 }=
$$
$$=\frac{K(1-T^2 \omega^2 )}{(1-T^2 \omega^2 )^2 +(2\xi T\omega )^2 }-j\frac{2K\xi T\omega }{(1-T^2 \omega^2 )^2 +(2\xi T\omega )^2 }$$

Весовая:

 $$
w(t)={\mathcal{L}}^{-1} \lbrace W(s)\rbrace ={\mathcal{L}}^{-1} \left\lbrace \frac{K}{T^2 s^2 +2\xi Ts+1}\right\rbrace ={\mathcal{L}}^{-1} \left\lbrace \frac{K/T^2 }{s^2 +2\frac{\xi }{T}s+\frac{1}{T^2 }}\right\rbrace =$$$$
={\mathcal{L}}^{-1} \left\lbrace \frac{K/T^2 }{{\left(s+\frac{\xi }{T}\right)}^2 +\frac{1-\xi^2 }{T^2 }}\right\rbrace ={\mathcal{L}}^{-1} \left\lbrace \frac{K/T^2 }{{\left(s+\frac{\xi }{T}\right)}^2 +{\left(\frac{\sqrt{1-\xi^2 }}{T}\right)}^2 }\right\rbrace =$$$$
=\frac{K/T^2 }{\frac{\sqrt{1-\xi^2 }}{T}}e^{-\frac{\xi }{T}t} \sin \left(\frac{\sqrt{1-\xi^2 }}{T}t\right)=\frac{K}{T\sqrt{1-\xi^2 }}e^{-\frac{\xi }{T}t} \sin \left(\frac{\sqrt{1-\xi^2 }}{T}t\right)
$$

Переходная:

 $$
h(t)={\mathcal{L}}^{-1} \left\lbrace \frac{W(s)}{s}\right\rbrace ={\mathcal{L}}^{-1} \left\lbrace \frac{K}{s(T^2 s^2 +2\xi Ts+1)}\right\rbrace ={\mathcal{L}}^{-1} \left\lbrace \frac{K}{s}-\frac{K(s+\frac{2\xi }{T})}{s^2 +\frac{2\xi }{T}s+\frac{1}{T^2 }}\right\rbrace =$$$$
={\mathcal{L}}^{-1} \left\lbrace \frac{K}{s}-K\cdot \frac{s+\frac{\xi }{T}}{(s+\frac{\xi }{T})^2 +{\left(\frac{\sqrt{1-\xi^2 }}{T}\right)}^2 }-\frac{K\xi }{T}\cdot \frac{1}{(s+\frac{\xi }{T})^2 +{\left(\frac{\sqrt{1-\xi^2 }}{T}\right)}^2 }\right\rbrace =$$$$
=K-Ke^{-\frac{\xi }{T}t} \left\lbrack \cos \left(\frac{\sqrt{1-\xi^2 }}{T}t\right)+\frac{\xi }{\sqrt{1-\xi^2 }}\sin \left(\frac{\sqrt{1-\xi^2 }}{T}t\right)\right\rbrack 
$$

АЧХ: 

$$A(\omega )=|W(j\omega )|=\sqrt{{\text{Re}}^2 (W)+{\text{Im}}^2 (W)}=\sqrt{\frac{K^2 }{(1-T^2 \omega^2 )^2 +(2\xi T\omega )^2 }}=\frac{K}{\sqrt{(1-T^2 \omega^2 )^2 +(2\xi T\omega )^2 }}$$

ФЧХ: $\varphi (\omega )=\arg W(j\omega )=\text{atan2}\left(-2\xi T\omega ,\,1-T^2 \omega^2 \right)$

ЛАЧХ: $L(\omega )=20\lg (A(\omega ))=20\lg K-10\lg \left\lbrack (1-T^2 \omega^2 )^2 +(2\xi T\omega )^2 \right\rbrack$

ЛФЧХ: $\varphi (\omega )=-\text{atan}\left(\frac{2\xi T\omega }{1-T^2 \omega^2 }\right)$


\textbf{Графики временных и частотных характеристик полученные в результате моделирования, их сравнение с построенными на основании аналитических выражений:}

Графики временных характеристик представлены на рис. \ref{time_dpt2}.
\begin{figure}[H]
    \centering
    \includegraphics[width=1\linewidth]{/home/ns/Desktop/ITMO-courses/Linear_automatic_control_systems/lab_5/code/images/plot_dpt2.png}
    \caption{Графики временных характеристик}
    \label{time_dpt2}
\end{figure}

Графики частотных характеристик представлены на рис. \ref{freq_dpt2}.
\begin{figure}[H]
    \centering
    \includegraphics[width=1\linewidth]{/home/ns/Desktop/ITMO-courses/Linear_automatic_control_systems/lab_5/code/images/freq_dpt2.png}
    \caption{Графики частотных характеристик}
    \label{freq_dpt2}
\end{figure}



\subsection{Конденсируй-умножай}

Дано уравнение конденсатора:

$I=C\frac{dU}{dt}$.

Возьмем из \textbf{Таблицы 2} инструкции по выполнению данной л/р значение, которое соответствует нашему варианту, для следующих величин:

• $C=354$ мкФ — ёмкость конденсатора.

Входом объекта считаем $I(t)$, а выходом $U(t)$.

\textbf{Найдём передаточную функцию и определим тип звена:}

$I=C\dot{U}$, Передаточная функция: $W(s)=\frac{K}{s}$, где $K=\frac{1}{C}=\frac{1}{3.54\cdot 10^{-4} }=2824.86$.

Типовое звено - \textit{идеальное интегрирующее.}

\textbf{Запишем аналитические выражения для временных и частотных характеристик:}

$$W(j\omega )=\frac{K}{j\omega }=-j\frac{K}{\omega }=0-j\frac{K}{\omega }$$

Весовая: $w(t)={\mathcal{L}}^{-1} \lbrace W(s)\rbrace ={\mathcal{L}}^{-1} \left\lbrace \frac{K}{s}\right\rbrace =K$

Переходная: $h(t)={\mathcal{L}}^{-1} \left\lbrace \frac{W(s)}{s}\right\rbrace ={\mathcal{L}}^{-1} \left\lbrace \frac{K}{s^2 }\right\rbrace =Kt$

АЧХ: $A(\omega )=|W(j\omega )|=\sqrt{{\text{Re}}^2 (W)+{\text{Im}}^2 (W)}=\sqrt{0^2 +{\left(\frac{K}{\omega }\right)}^2 }=\frac{K}{\omega }$

ФЧХ: $\varphi (\omega )=\arg W(j\omega )=\text{atan2}\left(-\frac{K}{\omega },\,0\right)=-\frac{\pi }{2}$

ЛАЧХ: $L(\omega )=20\lg (A(\omega ))=20\lg K-20\lg \omega$

ЛФЧХ: $\varphi (\omega )=-\frac{\pi }{2}$

\textbf{Графики временных и частотных характеристик полученные в результате моделирования, их сравнение с построенными на основании аналитических выражений:}

Графики временных характеристик представлены на рис. \ref{time_c}.
\begin{figure}[H]
    \centering
    \includegraphics[width=1\linewidth]{/home/ns/Desktop/ITMO-courses/Linear_automatic_control_systems/lab_5/code/images/plot_cap.png}
    \caption{Графики временных характеристик}
    \label{time_c}
\end{figure}

Графики частотных характеристик представлены на рис. \ref{freq_c}.
\begin{figure}[H]
    \centering
    \includegraphics[width=1\linewidth]{/home/ns/Desktop/ITMO-courses/Linear_automatic_control_systems/lab_5/code/images/freq_cap.png}
    \caption{Графики частотных характеристик}
    \label{freq_c}
\end{figure}

\subsection{Пружинка}

Даны уравнения пружинного маятника:

$F_{\text{упр}} =-k\cdot x$, $F=m\cdot a$

Возьмем из \textbf{Таблицы 3} инструкции по выполнению данной л/р значения, которые соответствуют нашему варианту, для следующих величин:

• $M=26$ кг — масса груза;

• $k=72$ Н/м — коэффициент жесткости пружины.

Входом объекта считаем $F_{ext} (t)$ (некую внешнюю силу, направленную соосно движению маятника), а выходом $x(t)$. Считаем, что маятник движется ортогонально силе тяжести (см. рисунок \ref{spring}).

\begin{figure}[H]
    \centering
    \includegraphics[width=0.3\linewidth]{/home/ns/Desktop/ITMO-courses/Linear_automatic_control_systems/lab_5/code/images/spring_pendulum.png}
    \caption{Схема}
    \label{spring}
\end{figure}

\textbf{Найдём передаточную функцию и определим тип звена:}

$$F=m\cdot a=m\ddot{x} ,~~m\ddot{x} =-kx+F_{ext} ,~~\frac{m}{k}\ddot{x} +x=\frac{F_{ext} }{k}$$

Передаточная функция: $W(s)=\frac{K}{T^2 s^2 +1}$, где $K=\frac{1}{k}=\frac{1}{72}=0.01389,~~T=\sqrt{\frac{m}{k}}=\sqrt{\frac{26}{72}}=\sqrt{0.3611}=0.6009$.

Типовое звено - \textit{консервативное}.

\textbf{Запишем аналитические выражения для временных и частотных характеристик:}

$$W(j\omega )=\frac{K}{1-T^2 \omega^2 }$$

Весовая: $w(t)={\mathcal{L}}^{-1} \lbrace W(s)\rbrace ={\mathcal{L}}^{-1} \left\lbrace \frac{K}{T^2 s^2 +1}\right\rbrace =\frac{K}{T}\sin \left(\frac{t}{T}\right)$

Переходная: $h(t)={\mathcal{L}}^{-1} \left\lbrace \frac{W(s)}{s}\right\rbrace ={\mathcal{L}}^{-1} \left\lbrace \frac{K}{s(T^2 s^2 +1)}\right\rbrace =K\left(1-\cos \left(\frac{t}{T}\right)\right)$

АЧХ: $A(\omega )=|W(j\omega )|=\sqrt{{\text{Re}}^2 (W)+{\text{Im}}^2 (W)}=\frac{K}{\left|1-T^2 \omega^2 \right|}$

ФЧХ: $\varphi (\omega )=\arg W(j\omega )=\left\lbrace \begin{array}{cc}
0, & \omega \le \frac{1}{T},\\
-\pi , & \omega >\frac{1}{T},
\end{array}
\right.$

ЛАЧХ: $L(\omega )=20\lg A(\omega )=20\lg K-20\lg \left(1-T^2 \omega^2 \right)$

ЛФЧХ: $\varphi (\omega )=\left\lbrace \begin{array}{cc}
0, & \omega \le \frac{1}{T},\\
-\pi , & \omega >\frac{1}{T},
\end{array}
\right.$

\textbf{Графики временных и частотных характеристик полученные в результате моделирования, их сравнение с построенными на основании аналитических выражений:}

Графики временных характеристик представлены на рис. \ref{time_s}.
Графики частотных характеристик представлены на рис. \ref{freq_s}.
\begin{figure}[H]
    \centering
    \includegraphics[width=1\linewidth]{/home/ns/Desktop/ITMO-courses/Linear_automatic_control_systems/lab_5/code/images/freq_spring.png}
    \caption{Графики частотных характеристик}
    \label{freq_s}
\end{figure}
\begin{figure}[H]
    \centering
    \includegraphics[width=1\linewidth]{/home/ns/Desktop/ITMO-courses/Linear_automatic_control_systems/lab_5/code/images/plot_spring.png}
    \caption{Графики временных характеристик}
    \label{time_s}
\end{figure}


\subsection{Что ты такое?}

На рисунке \ref{what} представлена принципиальная схема регулятора на операционном усилителе. Возьмем из \textbf{Таблицы 4} инструкции по выполнению данной л/р значения, которые соответствуют нашему варианту, для следующих величин:

• $R_{\text{1}} =2609$ Ом — сопротивление входного резистора;

• $R_{\text{2}} =13809$ Ом — сопротивление резистора отрицательной обратной связи;

• $C=266$  — емкость конденсатора отрицательной обратной связи.

Входом объекта считаем $U_{\text{ВХ}} (t)$, а выходом $U_{\text{ВЫХ}}$.

Помимо соответствия типовому звену определим тип регулятора.

\textit{В рамках этого задания разрешается не учитывать инверсию рассматриваемого регулятора.}

\begin{figure}[H]
    \centering
    \includegraphics[width=1\linewidth]{/home/ns/Desktop/ITMO-courses/Linear_automatic_control_systems/lab_5/code/images/principal_osc.png}
    \caption{Схема}
    \label{what}
\end{figure}

\textbf{Найдём передаточную функцию и определим тип звена:}

Составим систему на основе схемы:

$$I_1 =\frac{U_{\text{ВХ}} }{R_1 },~~\frac{U_{\text{ВХ}} }{R_1 }=-\frac{U_{\text{ВЫХ}} }{\left(R_2 +\frac{1}{Cs}\right)},~~W(s)=\frac{U_{\text{ВЫХ}} }{U_{\text{ВХ}} }=\frac{R_2 +\frac{1}{Cs}}{-R_1 }=\frac{\frac{1}{-R_1 C}(CR_2 s+1)}{s}$$

Передаточная функция: $W(s)=\frac{K(Ts+1)}{s}$, где $K=\frac{1}{-R_1 C}=\frac{1}{2609\cdot 266\cdot 10^{-6} }=1.441,$\\
$T=CR_2 =266\cdot 10^{-6} \cdot 13809=3.673$.

Типовое звено - \textit{изодромное. }Тип регулятора - \textit{пропорционально-интегральный (ПИ)}.

\textbf{Запишем аналитические выражения для временных и частотных характеристик:}

$$W(j\omega )=\frac{K(T(j\omega )+1)}{j\omega }=KT-\frac{K}{\omega }j$$

Весовая: $w(t)={\mathcal{L}}^{-1} \lbrace W(s)\rbrace ={\mathcal{L}}^{-1} \left\lbrace \frac{K(Ts+1)}{s}\right\rbrace =K(T\delta (t)+1)$

Переходная: $h(t)={\mathcal{L}}^{-1} \left\lbrace \frac{W(s)}{s}\right\rbrace ={\mathcal{L}}^{-1} \left\lbrace \frac{K(Ts+1)}{s^2 }\right\rbrace =K(T+t)$

АЧХ: 

$$A(\omega )=|W(j\omega )|=\sqrt{{\text{Re}}^2 (W)+{\text{Im}}^2 (W)}=\sqrt{(KT)^2 +{\left(\frac{K}{\omega }\right)}^2 }=K\sqrt{T^2 +\frac{1}{\omega^2 }}$$

ФЧХ: $\varphi (\omega )=\arg W(j\omega )=\text{atan2}\left(-\frac{K}{\omega },\,KT\right)$

ЛАЧХ: $L(\omega )=20\lg (A(\omega ))=20\lg K+10\lg \left(T^2 +\frac{1}{\omega^2 }\right)$

ЛФЧХ: $\varphi (\omega )=-\text{atan}\left(\frac{1}{T\omega }\right)$

\textbf{Графики временных и частотных характеристик полученные в результате моделирования, их сравнение с построенными на основании аналитических выражений:}

Графики временных характеристик представлены на рис. \ref{time_pi}.
\begin{figure}[H]
    \centering
    \includegraphics[width=1\linewidth]{/home/ns/Desktop/ITMO-courses/Linear_automatic_control_systems/lab_5/code/images/plot_pi.png}
    \caption{Графики временных характеристик}
    \label{time_pi}
\end{figure}

Графики частотных характеристик представлены на рис. \ref{freq_pi}.
\begin{figure}[H]
    \centering
    \includegraphics[width=1\linewidth]{/home/ns/Desktop/ITMO-courses/Linear_automatic_control_systems/lab_5/code/images/freq_pi.png}
    \caption{Графики частотных характеристик}
    \label{freq_pi}
\end{figure}

\subsection{Выводы}
В данной лабораторной работе мы исследовали несколько различных динамических объектов
и идентифицировали типовые звенья, которые им соответствуют. Для каждого объекта были получены
передаточные функции в стандартизированной форме через коэффициенты усиления $K$ и постоянные времени $T$,
записаны аналитические выражения для временных и частотных характеристик,
проведено моделирование и сравнение теоретических и экспериментальных результатов.\\
Во всех случаях наблюдается соответствие между теоретически рассчитанными характеристиками и результатами моделирования,
что подтверждает правильность идентификации типовых звеньев и полученных аналитических выражений.\\
Данная работа показывает, что типовые звенья довольно удобны, потому что при помощи них можно понятно описать самые разные
объекты, с которыми нужно будет работать на практике.

\newpage
\section{Приложение}

\matlabtitle{ДПТ}

\begin{matlabcode}
K = 2.615;
T = 0.0602;
simtime = 0.6;

out_h = sim("lab5_1_step.slx", simtime);

figure;
set(gcf, AutoResizeChildren="off", Position=[100 100 1200 1200]);

subplot(2,1,1);
hold on;

plot(out_h.y.Time, out_h.y.Data, "b", ...
     LineWidth=4, ...
     DisplayName="$h(t)$, simulation");

t = out_h.y.Time;
plot(t, K*(1 - exp(-t/T)), "r--", ...
     LineWidth=4, ...
     DisplayName="$h(t)$, analytical");

grid on;
legend(Location="best", FontSize=14, Interpreter="latex");
xlabel("$t, c$", Interpreter="latex");
ylabel("$h(t)$", Interpreter="latex");
set(gca, FontSize=14);
title("Step response", Interpreter="latex");
hold off;

subplot(2,1,2);
hold on;

out_w = sim("lab5_1_imp.slx", simtime);

plot(out_w.w.Time, out_w.w.Data, "b", ...
     LineWidth=4, ...
     DisplayName="$w(t)$, simulation");

t = out_w.w.Time;
plot(t, (K/T)*exp(-t/T), "r--", ...
     LineWidth=4, ...
     DisplayName="$w(t)$, analytical");

grid on;
legend(Location="best", FontSize=14, Interpreter="latex");
xlabel("$t, c$", Interpreter="latex");
ylabel("$w(t)$", Interpreter="latex");
set(gca, FontSize=14);
title("Impulse response", Interpreter="latex");
hold off;

exportgraphics(gcf, "images/plot_dpt.png", Resolution=200);
\end{matlabcode}


\begin{matlabcode}
K = 2.615;
T = 0.0602;

sys = tf(K, [T 1]);

w = linspace(0, 1000, 100000);

[mag, phase] = bode(sys, w);
mag = squeeze(mag);
phase = squeeze(phase);
phase = unwrap(phase * pi/180) * 180/pi;

A_analytic = K ./ sqrt((T*w).^2 + 1);
phi_analytic = -atan(T*w);

figure(Position=[100 100 1400 900]);
set(gcf, Color='w');

subplot(2,2,1);
plot(w, mag, ...
    Color='b', LineWidth=3, DisplayName='Simulation');
hold on;
plot(w, A_analytic, ...
    Color='r', LineStyle='--', LineWidth=3, DisplayName='Analytical');
grid on; box on;
xlabel('$\omega$, rad/s', Interpreter='latex');
ylabel('$A(\omega)$', Interpreter='latex');
title('AFR', Interpreter='latex');
legend(Location='best', Interpreter='latex', FontSize=12);

subplot(2,2,2);
plot(w, phase, ...
    Color='b', LineWidth=3, DisplayName='Simulation');
hold on;
plot(w, phi_analytic*180/pi, ...
    Color='r', LineStyle='--', LineWidth=3, DisplayName='Analytical');
grid on; box on;
xlabel('$\omega$, rad/s', Interpreter='latex');
ylabel('$\varphi(\omega), ^\circ$', Interpreter='latex');
title('PFR', Interpreter='latex');
legend(Location='best', Interpreter='latex', FontSize=12);


subplot(2,2,3);
semilogx(w, 20*log10(mag), ...
    Color='b', LineWidth=3, DisplayName='Simulation');
hold on;
semilogx(w, 20*log10(A_analytic), ...
    Color='r', LineStyle='--', LineWidth=3, DisplayName='Analytical');
grid on; box on;
xlabel('$\omega$, rad/s', Interpreter='latex');
ylabel('$L(\omega)$, dB', Interpreter='latex');
title('LAFR', Interpreter='latex');
legend(Location='best', Interpreter='latex', FontSize=12);

subplot(2,2,4);
semilogx(w, phase, ...
    Color='b', LineWidth=3, DisplayName='Simulation');
hold on;
semilogx(w, phi_analytic*180/pi, ...
    Color='r', LineStyle='--', LineWidth=3, DisplayName='Analytical');
grid on; box on;
xlabel('$\omega$, rad/s', Interpreter='latex');
ylabel('$\varphi(\omega), ^\circ$', Interpreter='latex');
title('LPFR', Interpreter='latex');
legend(Location='best', Interpreter='latex', FontSize=12);

set(findall(gcf, Type='axes'), FontSize=14);
exportgraphics(gcf, 'images/freq_dpt.png', Resolution=300);
\end{matlabcode}


\matlabtitle{ДПТ 2.0}

\begin{matlabcode}
K = 2.615;
T = 0.1207;
xi = 0.2668;
simtime = 3;

figure;
set(gcf, AutoResizeChildren="off", Position=[100 100 1200 1200]);

subplot(2,1,1);
hold on;

out_h = sim("lab5_2_step.slx", simtime);

plot(out_h.y.Time, out_h.y.Data, "b", ...
     LineWidth=4, ...
     DisplayName="$h(t)$, simulation");

t = out_h.y.Time;
wd = sqrt(1 - xi^2)/T;
h_analytic = K * (1 - exp(-xi*t/T) .* (cos(wd*t) + (xi/sqrt(1-xi^2))*sin(wd*t)));
plot(t, h_analytic, "r--", ...
     LineWidth=4, ...
     DisplayName="$h(t)$, analytical");

grid on;
legend(Location="best", FontSize=14, Interpreter="latex");
xlabel("$t, c$", Interpreter="latex");
ylabel("$h(t)$", Interpreter="latex");
set(gca, FontSize=14);
title("Step response", Interpreter="latex");
hold off;

subplot(2,1,2);
hold on;

out_w = sim("lab5_2_imp.slx", simtime);

plot(out_w.w.Time, out_w.w.Data, "b", ...
     LineWidth=4, ...
     DisplayName="$w(t)$, simulation");

t = out_w.w.Time;
wd = sqrt(1 - xi^2)/T;
w_analytic = (K/(T*sqrt(1-xi^2))) * exp(-xi*t/T) .* sin(wd*t);
plot(t, w_analytic, "r--", ...
     LineWidth=4, ...
     DisplayName="$w(t)$, analytical");

grid on;
legend(Location="best", FontSize=14, Interpreter="latex");
xlabel("$t, c$", Interpreter="latex");
ylabel("$w(t)$", Interpreter="latex");
set(gca, FontSize=14);
title("Impulse response", Interpreter="latex");
hold off;

exportgraphics(gcf, "images/plot_dpt2.png", Resolution=200);
\end{matlabcode}


\begin{matlabcode}
K = 2.615;
T = 0.1207;
xi = 0.2668;
sys = tf(K, [T^2, 2*xi*T, 1]);

w = linspace(0, 1000, 10000);

[mag, phase] = bode(sys, w);
mag = squeeze(mag); phase = squeeze(phase);
phase = unwrap(phase * pi/180) * 180/pi;

den = (1 - T^2*w.^2).^2 + (2*xi*T*w).^2;
A_analytic = K ./ sqrt(den);
phi_analytic = atan2(-2*xi*T*w, 1 - T^2*w.^2);

figure(Position=[100 100 1400 900]);
set(gcf, Color='w');

subplot(2,2,1);
plot(w, mag, ...
    Color='b', LineWidth=3, DisplayName='Simulation');
hold on;
plot(w, A_analytic, ...
    Color='r', LineStyle='--', LineWidth=3, DisplayName='Analytical');
grid on; box on;
xlabel('$\omega$, rad/s', Interpreter='latex');
ylabel('$A(\omega)$', Interpreter='latex');
title('AFR', Interpreter='latex');
legend(Location='best', Interpreter='latex', FontSize=12);
xlim([0, 100]);

subplot(2,2,2);
plot(w, phase, ...
    Color='b', LineWidth=3, DisplayName='Simulation');
hold on;
plot(w, phi_analytic*180/pi, ...
    Color='r', LineStyle='--', LineWidth=3, DisplayName='Analytical');
grid on; box on;
xlabel('$\omega$, rad/s', Interpreter='latex');
ylabel('$\varphi(\omega), ^\circ$', Interpreter='latex');
title('PFR', Interpreter='latex');
legend(Location='best', Interpreter='latex', FontSize=12);
xlim([0, 100]);

subplot(2,2,3);
semilogx(w, 20*log10(mag), ...
    Color='b', LineWidth=3, DisplayName='Simulation');
hold on;
semilogx(w, 20*log10(A_analytic), ...
    Color='r', LineStyle='--', LineWidth=3, DisplayName='Analytical');
grid on; box on;
xlabel('$\omega$, rad/s', Interpreter='latex');
ylabel('$L(\omega)$, dB', Interpreter='latex');
title('LAFR', Interpreter='latex');
legend(Location='best', Interpreter='latex', FontSize=12);

subplot(2,2,4);
semilogx(w, phase, ...
    Color='b', LineWidth=3, DisplayName='Simulation');
hold on;
semilogx(w, phi_analytic*180/pi, ...
    Color='r', LineStyle='--', LineWidth=3, DisplayName='Analytical');
grid on; box on;
xlabel('$\omega$, rad/s', Interpreter='latex');
ylabel('$\varphi(\omega), ^\circ$', Interpreter='latex');
title('LPFR', Interpreter='latex');
legend(Location='best', Interpreter='latex', FontSize=12);

set(findall(gcf, Type='axes'), FontSize=14);
exportgraphics(gcf, 'images/freq_dpt2.png', Resolution=300);
\end{matlabcode}


\matlabtitle{Конденсируй-умножай}

\begin{matlabcode}
K = 2824.86;
simtime = 1;

figure;
set(gcf, AutoResizeChildren="off", Position=[100 100 1200 1200]);

subplot(2,1,1);
hold on;

out_h = sim("lab5_3_step.slx", simtime);


plot(out_h.y.Time, out_h.y.Data, "b", ...
     LineWidth=4, ...
     DisplayName="$h(t)$, simulation");

t = out_h.y.Time;
plot(t, K*t, "r--", ...
     LineWidth=4, ...
     DisplayName="$h(t)$, analytical");

grid on;
legend(Location="best", FontSize=14, Interpreter="latex");
xlabel("$t, c$", Interpreter="latex");
ylabel("$h(t)$", Interpreter="latex");
set(gca, FontSize=14);
title("Step response", Interpreter="latex");
hold off;

subplot(2,1,2);
hold on;

out_w = sim("lab5_3_imp.slx", simtime);

plot(out_w.w.Time, out_w.w.Data, "b", ...
     LineWidth=4, ...
     DisplayName="$w(t)$, simulation");

t = out_w.w.Time;
plot(t, K*ones(size(t)), "r--", ...
     LineWidth=4, ...
     DisplayName="$w(t)$, analytical");

grid on;
legend(Location="best", FontSize=14, Interpreter="latex");
xlabel("$t, c$", Interpreter="latex");
ylabel("$w(t)$", Interpreter="latex");
set(gca, FontSize=14);
title("Impulse response", Interpreter="latex");
hold off;

exportgraphics(gcf, "images/plot_cap.png", Resolution=200);
\end{matlabcode}


\begin{matlabcode}
K = 2824.86;
sys = tf(K, [1 0]);
w = linspace(1e-3, 100, 100000);

[mag, phase] = bode(sys, w);
mag = squeeze(mag); phase = squeeze(phase);
phase = unwrap(phase * pi/180) * 180/pi;

A_analytic = K ./ w;
phi_analytic = -pi/2 * ones(size(w));

figure(Position=[100 100 1400 900]);
set(gcf, Color='w');

subplot(2,2,1);
plot(w, mag, ...
    Color='b', LineWidth=3, DisplayName='Simulation');
hold on;
plot(w, A_analytic, ...
    Color='r', LineStyle='--', LineWidth=3, DisplayName='Analytical');
grid on; box on;
xlabel('$\omega$, rad/s', Interpreter='latex');
ylabel('$A(\omega)$', Interpreter='latex');
title('AFR', Interpreter='latex');
legend(Location='best', Interpreter='latex', FontSize=12);
xlim([0, 0.1]);

subplot(2,2,2);
plot(w, phase, ...
    Color='b', LineWidth=3, DisplayName='Simulation');
hold on;
plot(w, phi_analytic*180/pi, ...
    Color='r', LineStyle='--', LineWidth=3, DisplayName='Analytical');
grid on; box on;
xlabel('$\omega$, rad/s', Interpreter='latex');
ylabel('$\varphi(\omega), ^\circ$', Interpreter='latex');
title('PFR', Interpreter='latex');
legend(Location='best', Interpreter='latex', FontSize=12);

subplot(2,2,3);
semilogx(w, 20*log10(mag), ...
    Color='b', LineWidth=3, DisplayName='Simulation');
hold on;
semilogx(w, 20*log10(A_analytic), ...
    Color='r', LineStyle='--', LineWidth=3, DisplayName='Analytical');
grid on; box on;
xlabel('$\omega$, rad/s', Interpreter='latex');
ylabel('$L(\omega)$, dB', Interpreter='latex');
title('LAFR', Interpreter='latex');
legend(Location='best', Interpreter='latex', FontSize=12);

subplot(2,2,4);
semilogx(w, phase, ...
    Color='b', LineWidth=3, DisplayName='Simulation');
hold on;
semilogx(w, phi_analytic*180/pi, ...
    Color='r', LineStyle='--', LineWidth=3, DisplayName='Analytical');
grid on; box on;
xlabel('$\omega$, rad/s', Interpreter='latex');
ylabel('$\varphi(\omega), ^\circ$', Interpreter='latex');
title('LPFR', Interpreter='latex');
legend(Location='best', Interpreter='latex', FontSize=12);

set(findall(gcf, Type='axes'), FontSize=14);
exportgraphics(gcf, 'images/freq_cap.png', Resolution=300);
\end{matlabcode}


\matlabtitle{Пружинка}

\begin{matlabcode}
K = 0.01389;
T = 0.6009;
simtime = 10;

figure;
set(gcf, AutoResizeChildren="off", Position=[100 100 1200 1200]);

subplot(2,1,1);
hold on;

out_h = sim("lab5_4_step.slx", simtime);

plot(out_h.y.Time, out_h.y.Data, "b", ...
     LineWidth=4, ...
     DisplayName="$h(t)$, simulation");

t = out_h.y.Time;
plot(t, K*(1 - cos(t/T)), "r--", ...
     LineWidth=4, ...
     DisplayName="$h(t)$, analytical");

grid on;
legend(Location="southoutside", FontSize=14, Interpreter="latex");
xlabel("$t, c$", Interpreter="latex");
ylabel("$h(t)$", Interpreter="latex");
set(gca, FontSize=10);
title("Step response", Interpreter="latex");
hold off;

subplot(2,1,2);
hold on;

out_w = sim("lab5_4_imp.slx", simtime);

plot(out_w.w.Time, out_w.w.Data, "b", ...
     LineWidth=4, ...
     DisplayName="$w(t)$, simulation");

t = out_w.w.Time;
plot(t, (K/T)*sin(t/T), "r--", ...
     LineWidth=4, ...
     DisplayName="$w(t)$, analytical");

grid on;
legend(Location="southoutside", FontSize=14, Interpreter="latex");
xlabel("$t, c$", Interpreter="latex");
ylabel("$w(t)$", Interpreter="latex");
set(gca, FontSize=10);
title("Impulse response", Interpreter="latex");
hold off;
exportgraphics(gcf, "images/plot_spring.png", Resolution=200);
\end{matlabcode}


\begin{matlabcode}
K = 0.01389;
T = 0.6009;
sys = tf(K, [T^2, 0, 1]);

w = linspace(0, 100, 10000);

[mag, phase] = bode(sys, w);
mag = squeeze(mag); phase = squeeze(phase);
phase = unwrap(phase * pi/180) * 180/pi;

A_analytic = K ./ abs(1 - T^2*w.^2);
phi_analytic = zeros(size(w));
phi_analytic(w > 1/T) = -pi;

figure(Position=[100 100 1400 900]);
set(gcf, Color='w');

subplot(2,2,1);
plot(w, mag, ...
    Color='b', LineWidth=3, DisplayName='Simulation');
hold on;
plot(w, A_analytic, ...
    Color='r', LineStyle='--', LineWidth=3, DisplayName='Analytical');
grid on; box on;
xlabel('$\omega$, rad/s', Interpreter='latex');
ylabel('$A(\omega)$', Interpreter='latex');
title('AFR', Interpreter='latex');
legend(Location='best', Interpreter='latex', FontSize=12);
xlim([0, 5]);

subplot(2,2,2);
plot(w, phase, ...
    Color='b', LineWidth=3, DisplayName='Simulation');
hold on;
plot(w, phi_analytic*180/pi, ...
    Color='r', LineStyle='--', LineWidth=3, DisplayName='Analytical');
grid on; box on;
xlabel('$\omega$, rad/s', Interpreter='latex');
ylabel('$\varphi(\omega), ^\circ$', Interpreter='latex');
title('PFR', Interpreter='latex');
legend(Location='best', Interpreter='latex', FontSize=12);
xlim([0, 5]);

subplot(2,2,3);
semilogx(w, 20*log10(mag), ...
    Color='b', LineWidth=3, DisplayName='Simulation');
hold on;
semilogx(w, 20*log10(A_analytic), ...
    Color='r', LineStyle='--', LineWidth=3, DisplayName='Analytical');
grid on; box on;
xlabel('$\omega$, rad/s', Interpreter='latex');
ylabel('$L(\omega)$, dB', Interpreter='latex');
title('LAFR', Interpreter='latex');
legend(Location='best', Interpreter='latex', FontSize=12);

subplot(2,2,4);
semilogx(w, phase, ...
    Color='b', LineWidth=3, DisplayName='Simulation');
hold on;
semilogx(w, phi_analytic*180/pi, ...
    Color='r', LineStyle='--', LineWidth=3, DisplayName='Analytical');
grid on; box on;
xlabel('$\omega$, rad/s', Interpreter='latex');
ylabel('$\varphi(\omega), ^\circ$', Interpreter='latex');
title('LPFR', Interpreter='latex');
legend(Location='best', Interpreter='latex', FontSize=12);

set(findall(gcf, Type='axes'), FontSize=14);
exportgraphics(gcf, 'images/freq_spring.png', Resolution=300);
\end{matlabcode}


\matlabtitle{Что ты такое?}

\begin{matlabcode}
K = 1.441;
T = 3.673;
simtime = 1;

figure;
set(gcf, AutoResizeChildren="off", Position=[100 100 1200 1200]);

subplot(2,1,1);
hold on;

out_h = sim("lab5_5_step.slx", simtime);

plot(out_h.y.Time, out_h.y.Data, "b", ...
     LineWidth=4, ...
     DisplayName="$h(t)$, simulation");

t = out_h.y.Time;
plot(t, K*(T + t), "r--", ...
     LineWidth=4, ...
     DisplayName="$h(t)$, analytical");

grid on;
legend(Location="best", FontSize=14, Interpreter="latex");
xlabel("$t, c$", Interpreter="latex");
ylabel("$h(t)$", Interpreter="latex");
set(gca, FontSize=14);
title("Step response", Interpreter="latex");
hold off;

subplot(2,1,2);
hold on;

out_w = sim("lab5_5_imp.slx", [ -0.1, simtime ]);

plot(out_w.w.Time, out_w.w.Data, "b", ...
     LineWidth=4, ...
     DisplayName="$w(t)$, simulation");

t = out_w.w.Time;
plot(t, K*ones(size(t)), "r--", ...
     LineWidth=4, ...
     DisplayName="$w(t)$, analytical");

grid on;
legend(Location="best", FontSize=14, Interpreter="latex");
xlabel("$t, c$", Interpreter="latex");
ylabel("$w(t)$", Interpreter="latex");
set(gca, FontSize=14);
title("Impulse response", Interpreter="latex");
ytickformat('%.2f')
ylim([1.4, 1.5]);
xlim([0, simtime]);
hold off;
exportgraphics(gcf, "images/plot_pi.png", Resolution=200);
\end{matlabcode}


\begin{matlabcode}
K = 1.441;
T = 3.673;
sys = tf(K*[T 1], [1 0]);

w = linspace(0, 1000, 10000);

[mag, phase] = bode(sys, w);    
mag = squeeze(mag); phase = squeeze(phase);
phase = unwrap(phase * pi/180) * 180/pi;

A_analytic = K * sqrt(T^2 + 1./w.^2);
phi_analytic = atan2(-1, T*w);

figure(Position=[100 100 1400 900]);
set(gcf, Color='w');

subplot(2,2,1);
plot(w, mag, ...
    Color='b', LineWidth=3, DisplayName='Simulation');
hold on;
plot(w, A_analytic, ...
    Color='r', LineStyle='--', LineWidth=3, DisplayName='Analytical');
grid on; box on;
xlabel('$\omega$, rad/s', Interpreter='latex');
ylabel('$A(\omega)$', Interpreter='latex');
title('AFR', Interpreter='latex');
legend(Location='best', Interpreter='latex', FontSize=12);
xlim([0, 10]);

subplot(2,2,2);
plot(w, phase, ...
    Color='b', LineWidth=3, DisplayName='Simulation');
hold on;
plot(w, phi_analytic*180/pi, ...
    Color='r', LineStyle='--', LineWidth=3, DisplayName='Analytical');
grid on; box on;
xlabel('$\omega$, rad/s', Interpreter='latex');
ylabel('$\varphi(\omega), ^\circ$', Interpreter='latex');
title('PFR', Interpreter='latex');
legend(Location='best', Interpreter='latex', FontSize=12);
xlim([0, 10]);

subplot(2,2,3);
semilogx(w, 20*log10(mag), ...
    Color='b', LineWidth=3, DisplayName='Simulation');
hold on;
semilogx(w, 20*log10(A_analytic), ...
    Color='r', LineStyle='--', LineWidth=3, DisplayName='Analytical');
grid on; box on;
xlabel('$\omega$, rad/s', Interpreter='latex');
ylabel('$L(\omega)$, dB', Interpreter='latex');
title('LAFR', Interpreter='latex');
legend(Location='best', Interpreter='latex', FontSize=12);

subplot(2,2,4);
semilogx(w, phase, ...
    Color='b', LineWidth=3, DisplayName='Simulation');
hold on;
semilogx(w, phi_analytic*180/pi, ...
    Color='r', LineStyle='--', LineWidth=3, DisplayName='Analytical');
grid on; box on;
xlabel('$\omega$, rad/s', Interpreter='latex');
ylabel('$\varphi(\omega), ^\circ$', Interpreter='latex');
title('LPFR', Interpreter='latex');
legend(Location='best', Interpreter='latex', FontSize=12);

set(findall(gcf, Type='axes'), FontSize=14);
exportgraphics(gcf, 'images/freq_pi.png', Resolution=300);
\end{matlabcode}
